\chapter{Objetivos y alcance}

\section{Objetivo general}
El objetivo general de esta tesis es diseñar, implementar y evaluar rigurosamente un defensor de ciberseguridad basado en aprendizaje por refuerzo, de carácter offline, que toma decisiones binarias sensibles al coste (\texttt{PERMIT} o \texttt{BLOCK}) sobre observaciones de flujos de red estandarizados. Este objetivo se persigue mediante un \textit{pipeline} experimental controlado que opera íntegramente sobre datos de flujo etiquetados y estáticos, priorizando el rigor metodológico y la reproducibilidad sobre cualquier pretensión de preparación para el despliegue.

El enfoque se fundamenta en la observación de que la detección de intrusiones implica costes inherentemente asimétricos: permitir un ataque conlleva consecuencias que son, en general, considerablemente más graves que las de generar una falsa alarma. Una formulación basada en recompensas hace explícita esta asimetría en el objetivo de aprendizaje, en lugar de dejarla implícita en un ajuste posterior de umbrales. La pregunta central que guía el trabajo es si un algoritmo de RL distribucional como QRDQN puede configurarse, entrenarse y evaluarse de forma que demuestre un comportamiento sensible al coste y medible sobre un benchmark público bien estudiado, manteniéndose a la vez comparable con sólidas líneas base supervisadas bajo las mismas condiciones controladas.

\section{Objetivos específicos}

\subsection{Representación del tráfico basada en flujos}
El primer objetivo específico es extraer y normalizar datos de tráfico de red en registros tabulares de flujos, garantizando que el formato de entrada sea compatible tanto con benchmarks públicos como CICIDS2017 como con capturas de laboratorio offline. Los flujos de red son resúmenes bidireccionales de intercambios de paquetes que comparten una quíntupla común, y su estructura tabular compacta y de anchura fija los hace muy adecuados como entradas para modelos de aprendizaje automático \parencite{Sperotto2010FlowIDS,Umer2017FlowIDS}. Alcanzar este objetivo requiere auditar las características disponibles en los ficheros CSV generados por CICFlowMeter, identificar qué campos contienen señal estadística legítima y cuáles representan artefactos de extracción o posibles vectores de fuga de información.

Este objetivo implica también definir un contrato de características estable que tienda un puente entre las publicaciones académicas en formato CSV y las exportaciones brutas de CICFlowMeter. Un contrato estable significa que el conjunto de características utilizado como entrada al modelo no varía entre experimentos, que los campos excluidos por razones de fuga se excluyen de forma consistente, y que cualquier persona que lea el código fuente puede verificar exactamente qué medidas llegan al modelo. Esta estabilidad es un requisito previo para todos los objetivos posteriores, ya que todos los demás componentes del pipeline dependen de un espacio de observación bien definido y auditable.

\subsection{Preprocesamiento canónico y esquema de características}
El segundo objetivo específico es definir y hacer cumplir un esquema canónico estricto de 76 características junto con una máscara de presencia/ausencia de 76 valores, lo que da lugar a un vector de observación estable de 152 dimensiones. Las 76 características canónicas se extraen de estadísticas de estilo CICFlowMeter que cubren la duración del flujo, conteos de paquetes y bytes, distribuciones de tiempos de llegada, conteos de flags TCP, estadísticas de longitud de paquetes, resúmenes de periodos activos e inactivos, y métricas de tasa de flujo. Los campos que podrían exponer información de identidad específica del conjunto de datos, como las direcciones IP de origen y destino, las marcas temporales absolutas, los \textit{Flow ID} y los números de puerto que funcionan como proxies directos del tipo de ataque, se excluyen explícitamente a nivel de esquema, en lugar de filtrarse en pasos de preprocesamiento ad hoc.

La máscara de presencia/ausencia de 76 valores registra si cada valor canónico fue medido directamente y está presente (codificado como 1) o ausente y sustituido por un marcador de posición por defecto (codificado como 0). Este diseño hace visibles las lagunas de datos tanto para el modelo como para el evaluador, evitando una situación en la que la red se entrena silenciosamente sobre ceros imputados sin ninguna indicación de que la medición subyacente no estaba disponible. Al pasar de CICIDS2017 a tráfico capturado en laboratorio, distintos pipelines de extracción pueden producir conjuntos de columnas diferentes; la máscara de presencia/ausencia garantiza que cualquier discrepancia de este tipo quede codificada en la observación en lugar de quedar oculta, de modo que el lector pueda distinguir un cero medido de uno imputado. Juntos, el esquema canónico y la máscara constituyen el contrato de observación del que depende cada componente posterior.

\subsection{Entorno de aprendizaje por refuerzo}
El tercer objetivo específico es construir una interfaz conjunto-de-datos-como-entorno compatible con Gymnasium \parencite{GymnasiumDocs} que presente los registros de flujo al agente como una secuencia de observaciones. El entorno envuelve un conjunto de datos etiquetado estático de modo que cada llamada a \texttt{step()} devuelve un registro de flujo como observación, junto con la recompensa calculada a partir de la acción anterior del agente y la etiqueta verdadera del último registro. Esta formulación permite que los bucles de entrenamiento de RL estándar, desarrollados para entornos interactivos, se apliquen directamente a conjuntos de datos de flujo offline sin modificar el algoritmo.

El entorno debe implementar una función de recompensa asimétrica y sensible al coste que penalice los falsos negativos más que los falsos positivos, reflejando los perfiles de riesgo en ciberseguridad en los que los ataques no detectados son operacionalmente más perjudiciales que las falsas alarmas \parencite{Lee2002CostSensitiveIDS,Cardenas2006IDSFramework}. Los pesos de recompensa específicos se tratan como parámetros experimentales configurables en lugar de constantes universales fijas, ya que la relación de costes adecuada depende del modelo de amenaza y del contexto operativo del escenario de despliegue. El entorno debe también admitir reinicios reproducibles, gestión de límites de episodios y particiones del conjunto de datos configurables, de modo que la misma interfaz pueda utilizarse tanto para el entrenamiento como para la evaluación bajo distintos modos de partición.

\subsection{Agente de decisión binaria basado en QRDQN}
El cuarto objetivo específico es configurar y entrenar un agente Quantile Regression Deep Q-Network \parencite{Dabney2018QRDQN}, aprovechando su enfoque distribucional para la estimación del valor, para navegar el espacio de acciones binario \texttt{PERMIT}/\texttt{BLOCK} a partir de las observaciones canónicas de 152 dimensiones. QRDQN extiende el DQN estándar estimando la distribución completa de los retornos esperados en lugar de solo la media, utilizando un conjunto de cuantiles aprendidos y una pérdida de regresión cuantílica. Esta representación distribucional resulta de interés en un contexto sensible al coste porque proporciona información más rica sobre la dispersión y la forma de los posibles resultados, no solo el retorno promedio, y en principio permite que la política sea más conservadora ante los peores casos.

La implementación práctica utiliza el módulo QRDQN de la biblioteca \texttt{sb3-contrib} \parencite{SB3ContribQRDQNDocs}, que proporciona una implementación estable y probada del algoritmo con arquitectura de red configurable, buffer de repetición, esquema de exploración e hiperparámetros de entrenamiento. El objetivo incluye seleccionar una configuración razonable por defecto, validar que el bucle de entrenamiento converge en el benchmark interno, y persistir todos los artefactos del modelo entrenado junto con la configuración y el estado de preprocesamiento utilizados para producirlos. La reproducibilidad exige que un modelo almacenado pueda cargarse y evaluarse de forma idéntica en una fecha posterior sin necesidad de reconstruir el pipeline de entrenamiento completo desde cero.

\subsection{Comparación con línea base supervisada}
El quinto objetivo específico es implementar una línea base supervisada robusta, utilizando principalmente Random Forest \parencite{LiuLang2019IDSSurvey}, para ofrecer una evaluación comparativa justa, con el mismo protocolo, del valor empírico del agente QRDQN. Random Forest se elige como línea base principal porque está bien establecido para la clasificación tabular, gestiona interacciones no lineales entre características sin un preprocesamiento extenso, es robusto ante características irrelevantes y tiende a obtener buenos resultados en benchmarks de NIDS basados en flujos \parencite{Apruzzese2022CrossEvaluationNIDS}. Utilizar el mismo esquema de características, el mismo protocolo de partición y las mismas métricas de evaluación tanto para el agente de RL como para la línea base es esencial; cualquier discrepancia en el preprocesamiento, la disponibilidad de características o las condiciones de evaluación comprometería la validez de la comparación.

Esta comparación cumple varios propósitos. En primer lugar, determina si la formulación QRDQN ofrece alguna ventaja empírica sobre un clasificador supervisado bien ajustado en condiciones controladas. En segundo lugar, proporciona una comprobación de cordura sobre la dificultad de la tarea en el benchmark, ya que un Random Forest que alcanza puntuaciones muy altas contextualiza lo que QRDQN logra. En tercer lugar, fundamenta cualquier discusión sobre el diseño de la función de recompensa sensible al coste en evidencia concreta. Si ambos modelos se evalúan utilizando las mismas métricas ponderadas por coste, la comparación revela si la función de recompensa de RL conduce realmente al agente hacia un punto de operación diferente al de un modelo supervisado entrenado con los mismos datos, o si ambos enfoques convergen a un comportamiento similar.

\subsection{Evaluación con control de fugas}
El sexto objetivo específico es diseñar pipelines de preprocesamiento y protocolos de validación que impidan estrictamente la fuga de información en todos los niveles de evaluación. Esto requiere la exclusión programática de campos identificativos, como las direcciones IP de origen y destino, las marcas temporales absolutas, los Flow ID y los proxies directos de puertos, en la fase de carga de datos antes de que comience cualquier entrenamiento de modelos \parencite{Arp2020DosDontsMLSecurity,Kaufman2011LeakageDataMining}. Estos campos se excluyen a nivel de esquema mediante el contrato canónico de características, en lugar de depender de un filtrado ad hoc que podría aplicarse de forma inconsistente en distintas ejecuciones o versiones del conjunto de datos.

Las transformaciones de escalado deben ajustarse exclusivamente sobre la partición de entrenamiento y luego aplicarse a las particiones de validación y prueba utilizando únicamente los parámetros aprendidos de los datos de entrenamiento. Ajustar un escalador sobre el conjunto de datos completo antes de la partición permite que la información distribucional del conjunto de prueba influya en la normalización de las características de entrenamiento, una forma sutil de fuga de información que puede inflar el rendimiento reportado en conjuntos de datos con una estructura temporal o de escenario marcada. El escalador ajustado se persiste como parte del artefacto de la ejecución para que pueda recargarse y aplicarse de forma idéntica durante la inferencia de la Fase 2 sobre tráfico de laboratorio, garantizando que se utilice la misma normalización para flujos fuera de distribución sin que ninguna información sobre su distribución llegue al estado del escalador.

\subsection{Análisis de escala de entrenamiento y eficiencia de datos}
El séptimo objetivo específico es evaluar el rendimiento tanto del agente de RL como de las líneas base supervisadas bajo distintas cantidades de datos de entrenamiento, identificando las características de eficiencia de datos de la formulación propuesta y evaluando con qué rapidez converge QRDQN en comparación con las alternativas supervisadas. Los métodos de RL profundo pueden ser muy demandantes en cuanto a muestras, y su dinámica de aprendizaje difiere sustancialmente de la de los métodos de árboles de decisión en conjunto. Comprender cuánto tráfico etiquetado requiere cada enfoque para alcanzar un rendimiento útil es prácticamente relevante: el tráfico de ataque etiquetado es costoso de recopilar, y los escenarios de despliegue reales pueden no tener acceso a los grandes conjuntos de datos equilibrados disponibles en los benchmarks académicos \parencite{Ring2019DatasetSurvey}.

Los experimentos de escala de entrenamiento utilizan subconjuntos de entrenamiento progresivamente más grandes extraídos de la misma partición, midiendo cómo cambian la precisión, el recall, el F1, la tasa de falsos positivos y la tasa de falsos negativos a medida que crece el presupuesto. La comparación entre QRDQN y Random Forest en cada nivel de presupuesto revela si la formulación de RL tiene un requisito mayor de muestras, si converge más lentamente, o si alcanza una meseta de rendimiento final diferente. Estos resultados se interpretan como evidencia interna sobre esta combinación específica de formulación y benchmark, no como afirmaciones generales sobre la eficiencia de datos de RL en todos los contextos de NIDS.

\subsection{Validación complementaria con tráfico de laboratorio}
El octavo objetivo específico es construir un pipeline de inferencia offline robusto capaz de ingerir tráfico capturado por el operador en un entorno de laboratorio doméstico cerrado y no adversarial, y producir predicciones a partir de un modelo entrenado en CICIDS2017. Este nivel de validación ofrece una indicación parcial, de validez externa limitada, de si el esquema canónico de características y el modelo entrenado se comportan de forma razonable sobre flujos extraídos de un entorno de red distinto del benchmark, potencialmente utilizando una versión distinta de CICFlowMeter o una configuración de captura diferente. El cambio de dominio entre la distribución de entrenamiento y la distribución de inferencia es un reto fundamental para los modelos de NIDS y rara vez se evalúa en estudios que solo evalúan sobre un único conjunto de datos público \parencite{Apruzzese2022CrossEvaluationNIDS,Layeghy2023CrossDomainNIDS}.

El pipeline de inferencia aplica recorte opcional por percentiles y z-score para gestionar los valores de características brutas fuera de distribución que caen fuera del rango observado durante el entrenamiento en CICIDS2017. El escalador ajustado de la Fase 1 se carga y aplica sin reajuste, preservando la integridad de la separación entre entrenamiento y prueba. La máscara de presencia/ausencia en el vector de observación codifica explícitamente cualquier característica que no haya podido calcularse a partir de la captura de laboratorio, haciendo visible el caso de datos incompletos en la entrada de predicción en lugar de imputarlo silenciosamente. El resultado esperado de este objetivo es un conjunto de artefactos de predicción y métricas que permitan comparar directamente los resultados del tráfico de laboratorio con los de la validación de la Fase 1, proporcionando evidencia concreta sobre cómo varía el rendimiento bajo un cambio de dominio.

\section{Alcance del trabajo}
El trabajo está delimitado por los requisitos de una evaluación offline de un pipeline de aprendizaje automático tabular. Comprende la ingesta de registros de flujo en formato CSV tanto de CICIDS2017 como de capturas de laboratorio, la aplicación de escalado estándar y recorte opcional de valores atípicos, la ejecución de bucles de entrenamiento mediante el framework Stable Baselines3 (\texttt{sb3-contrib}), y la generación de artefactos de ejecución estructurados y reproducibles que persisten modelos, métricas, escaladores y ficheros de configuración junto a cada ejecución de entrenamiento. La experimentación interna se realiza íntegramente sobre el conjunto de datos CICIDS2017, utilizando particiones explícitas entre entrenamiento y prueba y una escalera de validación multietapa para evaluar el comportamiento del modelo bajo condiciones progresivamente más estrictas antes de probarlo con tráfico capturado externamente.

El alcance incluye explícitamente la escalera de validación multietapa descrita en los objetivos de evaluación. Esta escalera cubre particiones aleatorias estratificadas para comprobaciones básicas de aprendizaje, una comprobación de etiquetas barajadas para verificar que el rendimiento colapsa cuando las etiquetas no contienen señal real, una partición por día o por CSV para comprobar la generalización dentro de CICIDS2017, una evaluación \textit{leave-one-CSV-out} para comprobar la sensibilidad a ficheros de captura específicos, y la inferencia offline sobre flujos capturados por el operador en un laboratorio doméstico cerrado. Cada nivel aporta evidencia diferenciada y forma parte del alcance experimental acordado.

\section{No objetivos}
Esta tesis excluye explícitamente el desarrollo de monitorización de red en línea o la interceptación de paquetes en tiempo real. El sistema no se integrará con cortafuegos, reglas de \texttt{iptables}, controladores de Software-Defined Networking (SDN), ni sistemas de prevención de intrusiones (IPS) activos. Estas exclusiones existen porque el despliegue en producción requeriría un diseño de sistema sustancialmente diferente, una validación de seguridad y una metodología de pruebas operacionales que quedan fuera del alcance de un estudio académico controlado. La tesis está diseñada como un prototipo experimental, no como un producto operativo destinado a su uso en producción.

La tesis tampoco pretende resolver el problema más amplio de las operaciones cibernéticas adversariales multiagente o la interrupción de ataques con dependencia de secuencias, donde el estado del entorno cambia dinámicamente en respuesta a las acciones del defensor. En la formulación actual, cada registro de flujo se clasifica de forma independiente, y el conjunto de datos no reacciona a las decisiones del agente. Esta es una limitación conocida del enfoque \textit{conjunto-de-datos-como-entorno}, e implica que el sistema no puede modelar la adaptación adversarial, la persistencia del atacante ni las cadenas de ataques multietapa que abarcan muchas decisiones individuales de flujo. Por último, el objetivo no es demostrar que QRDQN represente el estado del arte definitivo para todas las tareas de NIDS, sino evaluar su comportamiento bajo un marco experimental cuidadosamente controlado, libre de fugas y reproducible, y compararlo honestamente con líneas base supervisadas.

\section{Resultados esperados}
El resultado esperado de esta tesis es un pipeline de inferencia offline completamente funcional, documentado y reproducible, respaldado por un análisis empírico detallado. El proyecto producirá artefactos persistentes para todas las ejecuciones de entrenamiento y validación, incluyendo ficheros de modelo entrenado, escaladores ajustados, informes de métricas, salidas de validación y registros de configuración que especifican completamente las condiciones bajo las cuales se produjo cada resultado. Esta disciplina en la gestión de artefactos forma parte de la contribución metodológica: un resultado que no puede reproducirse a partir de su rastro de artefactos proporciona una evidencia más débil que uno que sí puede.

En cuanto al contenido empírico, el proyecto producirá mediciones de precisión, recall, F1, tasa de falsos positivos y tasa de falsos negativos tanto para el agente QRDQN como para la línea base de Random Forest en todos los niveles de validación. Los experimentos de escala de entrenamiento producirán curvas de aprendizaje que caracterizarán cómo evoluciona el rendimiento de cada modelo a medida que crece el presupuesto de entrenamiento. Los resultados de inferencia de la Fase 2 caracterizarán cómo varía el rendimiento cuando el modelo se enfrenta a tráfico capturado por el operador en un laboratorio doméstico cerrado y no adversarial, en lugar de flujos de CICIDS2017 reservados para la prueba. En conjunto, estos resultados cuantificarán los límites de generalización de los modelos de RL entrenados en benchmarks cuando se someten a particiones temporales estrictas, evaluación \textit{leave-one-CSV-out} y tráfico de laboratorio doméstico generado por el operador, y proporcionarán una evaluación empírica honesta de si la formulación de RL sensible al coste ofrece ventajas medibles sobre una línea base supervisada bien equiparable.
